% Solutions of Book 4, chapter 6 (rust-ownership).
\section*{Chapter \ref{ch:b2:rust-ownership} --- Rust: Ownership and Borrowing}

\begin{solution}{exo:b2:rust-ownership:1}
\omoutput{code/bachelor-2/06-rust-ownership/outputs/run_compile.txt}{Each snippet built alone.}
a: two shared \omterm{def:b2:rust-ownership:borrow}{borrows} in use together, allowed. b: \code{first} \omterm{def:b2:rust-ownership:borrow}{borrows} \code{v}, \code{push} needs an exclusive
\omterm{def:b2:rust-ownership:borrow}{borrow} while \code{first} is still to be printed: refused (and rightly: the push may reallocate). c: \code{i32} is a
\omterm{def:b2:rust-ownership:copy-type}{copy type}, so \code{n} stays usable. d: \code{s} moved into \code{calc}, then used: refused. e: the exclusive \omterm{def:b2:rust-ownership:borrow}{borrow}
\code{r} is last used on the line before \code{t} is created, so the two are never in use together: allowed.
\end{solution}

\begin{solution}{exo:b2:rust-ownership:2}
\code{shortest} is a shared \omterm{def:b2:rust-ownership:borrow}{borrow} of an element of \code{names}, in use until the \code{println!}; \code{push}
needs an exclusive \omterm{def:b2:rust-ownership:borrow}{borrow} of \code{names} in between. The push may reallocate the vector's buffer, and
\code{shortest} would dangle. Fix 1, shorten: push before computing \code{shortest}, or print first and push after.
Fix 2, clone: \code{let shortest = names.iter().min_by_key(|n| n.len()).unwrap().clone();}, an owned copy of one
short string. A third, restructure: keep the index of the shortest name and read \code{names[i]} after the push.
\end{solution}

\begin{solution}{exo:b2:rust-ownership:3}
\omcode{code/bachelor-2/06-rust-ownership/exercises/first_word.rs}{2}{8}{The first word, borrowed}
One reference parameter: by the elision rules the result \omterm{def:b2:rust-ownership:borrow}{borrows} from \code{s}. \code{find} returns the byte index at
which a character begins (here the whitespace), so \code{&s[..i]} ends on a boundary. The tests cover Café,
a tab, a single word and a leading space, which gives the empty string.
\end{solution}

\begin{solution}{exo:b2:rust-ownership:4}
\omoutput{code/bachelor-2/06-rust-ownership/outputs/exo_clones.txt}{The count.}
\code{String::from} 1, \code{a.clone()} 2, \code{Vec::new()} none, the first \code{push} 3 (room for four),
\code{a.clone()} 4, \code{v[0].clone()} 5; \code{+ " Lane"} appends to that clone, whose buffer of exactly 4 bytes is
enlarged: the reallocation. \code{&c} allocates nothing; \code{to_uppercase} builds a new string: 6. \omterm{def:b2:rust-ownership:move}{Moves}
(\code{v.push(b)}, the result of \code{+}) allocate nothing.
\end{solution}

\begin{solution}{exo:b2:rust-ownership:5}
\omoutput{code/bachelor-2/06-rust-ownership/outputs/exo_drop.txt}{The ten lines.}
\code{let _ = Noisy("c");} binds nothing: the pattern \code{_} does not take \omterm{def:b2:rust-ownership:ownership}{ownership}, so the temporary is dropped at
the end of the statement, first of all. \code{b} \omterm{def:b2:rust-ownership:move}{moves} into \code{consume} and is dropped at the end of that function.
\code{f} at the end of its block. \code{drop(v)} \omterm{def:b2:rust-ownership:drop}{drops} the vector, whose elements are dropped in order, \code{d} then
\code{e}. The second \code{_a} shadows the first without dropping it (a name starting with \code{_} still owns); at
the end of \code{main} the locals are dropped in reverse order of declaration: \code{a2}, then \code{a}.
\end{solution}

\begin{solution}{exo:b2:rust-ownership:6}
\code{line} is a \code{String} owned by \code{joined} and dropped when it returns; \code{&line} would point to a
freed buffer. The fix: return the owned value, \code{fn joined(names: &[String]) -> String}, with \code{line} as the
last expression; the buffer \omterm{def:b2:rust-ownership:move}{moves} to the caller. The C version, returning a \code{char *} to a \code{malloc}ed
buffer that it freed, or to a local array, compiles; the caller reads freed memory or a dead stack frame, which
usually still holds the right bytes until the next call overwrites them, so the bug shows later and elsewhere.
\end{solution}

\begin{solution}{exo:b2:rust-ownership:7}
\omcode{code/bachelor-2/06-rust-ownership/exercises/split.rs}{2}{20}{Two halves, two exclusive \omterm{def:b2:rust-ownership:borrow}{borrows}}
The \omterm{def:b2:rust-ownership:checker}{borrow checker} reasons on whole values: \code{&mut v[..m]} \omterm{def:b2:rust-ownership:borrow}{borrows} \code{v}, and so does \code{&mut v[m..]},
whatever the ranges. \code{split_at_mut} is a library function whose body uses \code{unsafe} code to build the two
\omterm{def:b2:rust-ownership:slice}{slices} from one pointer; its signature returns two \code{&mut [T]} that \omterm{def:b2:rust-ownership:borrow}{borrow} \code{v}, and the library guarantees,
by construction, that they do not overlap. \code{tests/ownership} checks both functions, including an odd length,
equal largest elements and an empty \omterm{def:b2:rust-ownership:slice}{slice}.
\end{solution}

\begin{solution}{exo:b2:rust-ownership:8}
The signature promises that the result is valid wherever \code{x} is, and says nothing of \code{y}, whose lifetime
is unrelated. A caller may thus keep the result after \code{y} is dropped:
\code{let r; { let t = String::from("tmp"); r = f(&a, &t); } println!("{r}");} compiles. If \code{f} returned
\code{y}, \code{r} would dangle; the compiler refuses the body instead. Giving both the same lifetime,
\code{fn f<'a>(x: &'a str, y: &'a str) -> &'a str}, allows returning either; the caller loses the call above, which is
then refused because \code{t} does not live long enough.
\end{solution}

\begin{solution}{exo:b2:rust-ownership:9}
\omcode{code/bachelor-2/06-rust-ownership/exercises/reverse_words.rs}{2}{10}{Words reversed in place}
Let the text be $w_1\,w_2 \dots w_k$, words separated by spaces. Reversing all bytes gives
$\overline{w_k}\,\dots\,\overline{w_1}$, each word's bytes in reverse order; reversing each word gives back
$w_k \dots w_1$. Each multi-byte character lies inside one word, so its bytes are reversed twice and recover their
order; spaces are single bytes. The final bytes are therefore the concatenation of valid encoded characters, valid
UTF-8, and \code{from_utf8} succeeds without copying. \code{into_bytes} and \code{from_utf8} reuse the buffer: the
test compares the addresses.
\end{solution}

\begin{solution}{exo:b2:rust-ownership:10}
\omcodes{code/bachelor-2/06-rust-ownership/exercises/stack.rs}{1-30,49-57}{A stack of boxed nodes}
\code{self.head} cannot be moved out of a borrowed \code{self}, since \code{self} would be left without a head;
\code{take} \omterm{def:b2:rust-ownership:move}{moves} the value out and puts \code{None} in its place, a legal exclusive use of the field. The default \omterm{def:b2:rust-ownership:drop}{drop}
of a \code{Box<Node>} \omterm{def:b2:rust-ownership:drop}{drops} its \code{next} field, which \omterm{def:b2:rust-ownership:drop}{drops} the next box, and so on: one nested call per node, and
a million nested calls overflow the 8\,MB stack of the main thread. The \code{Drop} above detaches each node before
dropping it, so each node is dropped with no successor. The test pushes a million values and \omterm{def:b2:rust-ownership:drop}{drops} the stack, and
compares 10\,000 random operations with a \code{Vec}.
\end{solution}

\begin{solution}{exo:b2:rust-ownership:11}
(a) A loop over a vector holds an iterator, which is a shared \omterm{def:b2:rust-ownership:borrow}{borrow} of the vector (or an exclusive one, for
\code{iter_mut}), in use from the start of the loop to its last iteration. \code{push} takes \code{&mut self}, an
exclusive \omterm{def:b2:rust-ownership:borrow}{borrow}; a call inside the body would make it in use together with the iterator's \omterm{def:b2:rust-ownership:borrow}{borrow}, which the
\omterm{def:b2:rust-ownership:borrow}{borrowing rule} forbids. Without \code{unsafe}, there is no other path to the vector. (b)
\omoutput[2-12]{code/bachelor-2/06-rust-ownership/outputs/invalidation.txt}{The same loop, in Rust.}
The message names exactly the two \omterm{def:b2:rust-ownership:borrow}{borrows}. (c) Collect the renewals in a new vector, with \code{iter},
\code{filter} and \code{map} and a final \code{collect}, then append them after the loop with \code{due.extend}.
That was the first fix of the C++ chapter; it is now the one the compiler accepts. The index loop works too, since
\code{due[i]} \omterm{def:b2:rust-ownership:borrow}{borrows} only for the duration of the expression.
\end{solution}

\begin{solution}{exo:b2:rust-ownership:12}
\code{seen.clone()} copies the vector and every string in it, at every step. With $n = 2m$ names, the first $m$
distinct, step $i$ clones $k_i = \min(i, m)$ strings, which costs $k_i + 1$ allocations when $k_i > 0$ (the strings and
the vector's buffer). Adding the $m$ clones pushed and the $g = 1 + \lceil \log_2 (m / 4) \rceil$ buffers of
\code{seen}, the total is $\sum_{k=1}^{m-1} (k+1) + m(m+1) + m + g = \frac{3}{2}m^2 + \frac{5}{2}m - 1 + g$: 3879,
15\,255 and 60\,506 for $n = 100$, 200 and 400, the printed values, and $\Theta(n^2)$. The clone was never needed:
\code{seen.contains(n)} \omterm{def:b2:rust-ownership:borrow}{borrows} \code{seen} only for the call, so \code{seen.push} after it is legal.
\omcode{code/bachelor-2/06-rust-ownership/exercises/distinct.rs}{4}{18}{Without the clones, then with a hash set}
\code{distinct} still compares each name with the list, $O(nm)$ comparisons but only $m$ allocations;
\code{distinct_fast} keeps borrowed \code{&str}s in a hash set and is $O(n)$ expected. The test checks that both
give the same names in the order of first appearance.
\end{solution}

\begin{solution}{pb:b2:rust-ownership:1}
\textbf{1.} This counter is case-sensitive: \code{The} and \code{the} are two words. The previous chapter lowered every
word first.

\textbf{2.} \code{first}, the pointer to the counter of \code{The}, at line 18. The block was freed by \code{grow},
called by \code{dict_count} at line 17 when inserting a new word into a half-full table.

\textbf{3.} It reads freed memory only if an insertion grew the table between line 16 and line 18. An input whose
distinct words never reach half the capacity, say one line \code{The The}, never grows the table of 8 slots.

\textbf{4.} The table: \code{dict_count} stored a \code{strdup} copy, and \code{dict_free} frees it. \code{best} only
\omterm{def:b2:rust-ownership:borrow}{borrows} it. \code{main} frees it, then \code{dict_free} frees it again.

\textbf{5.} Into \code{line}, the buffer of \code{getline}, which is overwritten by every new line and freed after the
loop. Every stored word points into that one buffer, so after the next line is read the old words have changed: the
comparisons of later lookups compare with the wrong text, and counts go to the wrong words.

\textbf{6.} Version 1 increments, or here only reads, a freed slot that \code{malloc} may have given to another
use: silent corruption. Version 2 corrupts the allocator's own lists, a crash later or an exploitable state. Version
3 prints wrong counts, or garbage, with no crash at all. Version 3 is the most dangerous: it is wrong on every input
and may look plausible.

\textbf{7.} Each \code{String} is owned by its slot, the slots by the vector, the vector by the \code{Dict}. When the
\code{Dict} is dropped, the vector \omterm{def:b2:rust-ownership:drop}{drops} each \code{Option}, which \omterm{def:b2:rust-ownership:drop}{drops} each \code{String}, freeing its buffer, then
the vector frees its own buffer.

\textbf{8.} \code{mem::replace} \omterm{def:b2:rust-ownership:move}{moves} the old vector out of \code{self.slots}, leaving the new empty one, so the
function now owns the old one; \code{into_iter} \omterm{def:b2:rust-ownership:move}{moves} each pair out of it, and moving a \code{String} copies its three
words, not its bytes. Reading \code{self.slots} while writing a new vector stored in \code{self.slots} would need two
\omterm{def:b2:rust-ownership:borrow}{borrows} of the same field, one exclusive; after the replace they are two different vectors.

\textbf{9.} The result \omterm{def:b2:rust-ownership:borrow}{borrows} from \code{self} (a method with \code{&mut self}): the \code{Dict} stays exclusively
borrowed for as long as the returned counter is used.

\textbf{10.} The \code{Dict} is dropped automatically at the end of \code{main}, and the \omterm{def:b2:rust-ownership:drop}{drop} of its fields frees
everything: the compiler wrote \code{dict_free}.

\textbf{11.} \code{count} copies the word (\code{to_string}) when it creates a slot, and only compares it otherwise;
the \omterm{def:b2:rust-ownership:borrow}{borrow} \code{&str} is needed only during the call, which ends before the line is dropped.

\textbf{12.} The 34 \code{String} buffers, one by one in the order of the slots, then the slots' buffer; the three
words of \code{d} on the stack disappear with the frame of \code{main}.

\textbf{13.} \code{first} is an exclusive \omterm{def:b2:rust-ownership:borrow}{borrow} of \code{d}, in use until line 12; the second \code{count} needs
another. The comment's rule, valid until the next insertion, is exactly this \omterm{def:b2:rust-ownership:borrow}{borrow}: the Rust signature makes it a
rule the compiler checks.

\textbf{14.} Only the \omterm{def:b2:rust-ownership:ownership}{owner} frees: a \code{String} is freed when its \omterm{def:b2:rust-ownership:ownership}{owner} \omterm{def:b2:rust-ownership:drop}{drops} it. Through \code{&String} nothing
can take \omterm{def:b2:rust-ownership:ownership}{ownership}, so nothing can free it: E0507.

\textbf{15.} The map stores \code{&str} borrowed from \code{line}, which is dropped at the end of each iteration while
the map lives on: E0597, the lifetime of the \omterm{def:b2:rust-ownership:borrow}{borrow} exceeds that of the \omterm{def:b2:rust-ownership:ownership}{owner}.

\textbf{16.} 1: finish with \code{first} before inserting, or call \code{count("The")} again when needed. 2: clone
the word if an owned copy is needed: \code{let mine = best.clone();}. 3: store owned words,
\code{HashMap<String, u32>}, inserting \code{w.to_string()} when a word is new, as \code{Dict} does.

\textbf{17.} Version 3 disappears with \code{std::string} keys, which own copies. Versions 1 and 2 remain: a pointer
or reference into an \code{unordered_map} is stable across rehashes (the nodes do not \omterm{def:b2:rust-ownership:move}{move}), so version 1 would even
be correct there, but with a \code{vector} of slots it reads freed memory; and \code{delete} on a pointer obtained
from \code{c_str()} compiles. C++ checks neither.

\textbf{18.} The comment states a rule for humans; the signature
\code{fn count(&mut self, word: &str) -> &mut u32} states the same rule in a form the compiler enforces at every call.

\textbf{19.} Both print 100\,000 different words and the same count for \code{w7}: same generator, same counts.

\textbf{20.} On the printed run, both means are about a third of a second and lie within 5\,\% of each other, while
the three runs of one program differ by up to about 25\,\%: the difference between the means is smaller than the
spread.

\textbf{21.} The index is masked with \code{cap - 1}, and the check is one comparison with the length, predicted
correctly every time; the probing loop is dominated by hashing, string comparison and cache misses.

\textbf{22.} Formatting the word (\code{snprintf} against \code{write!}), the hash written as a fold, \code{strcmp}
against a \omterm{def:b2:rust-ownership:slice}{slice} comparison that first compares lengths, and the allocators' handling of 100\,000 small blocks.

\textbf{23.} Many runs on a quiet machine, the order of the runs alternated, the median and its spread reported, and
both programs compiled by the same back end (\code{gcc} and \code{rustc} optimise differently).

\textbf{24.} \emph{Three memory errors refused at compile time, same speed within a measured margin}: the use after
free across an insertion (E0499), the double free of a borrowed word (E0507) and the words borrowed from a dead
buffer (E0597) are all compile-time errors in Rust; and on the laptop the two versions run within 5\,\% of each other,
inside the run-to-run spread.
\end{solution}
